\section{Introduction}

AI coding agents are remarkably capable at code generation: controlled
studies of GitHub Copilot report productivity gains of 21--55\% across
programming tasks~\cite{peng2023copilot}, and the broader ecosystem can
now scaffold entire applications from natural-language descriptions.
Yet the dominant workflow pattern---what
Karpathy~\cite{karpathy2025vibe} termed ``vibe coding''---prioritizes
speed of implementation over deliberate preparation. The developer
describes an intent, the agent produces code, and misalignments are
resolved through iterative correction. This creates a systematic
alignment problem: agents operating without sufficient context produce
code that requires extensive rework. Veracode reports that 45\% of
AI-generated code contains security flaws~\cite{veracode2025ai}, and
practitioners consistently observe that debugging misaligned agent
output can consume substantial development
time~\cite{huntley2025ralph, yegge2025beads}. The bottleneck in agentic
coding is not code generation but \emph{alignment} between developer
intent and agent output---specification compliance, architectural
fidelity, and low corrective-commit ratios.

We argue that this alignment problem is fundamentally a \emph{preparation}
problem. Drawing on the culinary concept of \emph{mise en place} (MEP)---a
French term meaning ``everything in its place''---we propose a
preparation-first methodology for agentic coding. In professional kitchens,
thorough preparation enables fluid execution: every ingredient is
measured and every tool positioned so that once cooking begins, the
chef's hands never pause to search. The same principle applies to agent
orchestration. When domain expertise, design intent, and task boundaries
are externalized into structured artifacts \emph{before} agents begin
writing code, the resulting implementation is more aligned, more
coherent, and less costly to verify. MEP formalizes this preparation
into three steps: contextual grounding (tacit knowledge captured in
machine-legible documents), collaborative specification (human-agent
dialogue producing detailed design artifacts), and task decomposition
(specifications converted into dependency-aware work records). We
position MEP relative to spec-driven development, prompt engineering,
and iterative AI coding in Section~\ref{sec:related}.

This paper makes four contributions:
\begin{enumerate}
  \item A three-phase MEP methodology for agentic coding, grounded
    in backward design~\cite{wiggins1998understanding} and tacit knowledge
    externalization~\cite{polanyi1966tacit}.
  \item A case study applying this methodology in a competitive hackathon,
    with quantitative data on preparation artifacts and implementation
    outcomes.
  \item The concept of \emph{context fluency} as an emerging developer
    skill---the ability to create rich, structured context that AI agents can
    act on.
  \item A research agenda with five open questions for empirical validation
    of preparation-phase methodologies.
\end{enumerate}
